\documentclass[a4paper,11pt]{article}

% ---------- Packages ----------
\usepackage[margin=1in, left=0.7in, right=0.7in, top=0.6in, bottom=0.6in]{geometry}
\usepackage{lmodern}
\usepackage[T1]{fontenc}
\usepackage[utf8]{inputenc}
\usepackage{hyperref}
\usepackage{xcolor}
\usepackage{titlesec}
\usepackage{enumitem}
\usepackage{tabularx}
\usepackage{multicol}
\usepackage{parskip}

% ---------- Colors ----------
\definecolor{linkblue}{RGB}{0,0,128}
\definecolor{ruleblue}{RGB}{52,101,164}

% ---------- Hyperref setup ----------
\hypersetup{
    colorlinks=true,
    linkcolor=linkblue,
    urlcolor=linkblue,
    citecolor=linkblue
}

% ---------- Section formatting ----------
\titleformat{\section}
  {\large\bfseries}
  {}
  {0em}
  {}
  [\vspace{0.15em}\hrule\vspace{0.3em}]

\titlespacing*{\section}{0pt}{2.4em}{0.2em}

% ---------- Custom commands ----------
\newcommand{\cventry}[3]{%
  \noindent
  \begin{tabularx}{\textwidth}{@{}>{\small\raggedright\arraybackslash}p{2.2cm}|@{\hspace{0.4cm}}X@{}}
    #1 & \textbf{#2} \\
    & #3
  \end{tabularx}\vspace{0.3em}
}

\newcommand{\cventrynobar}[3]{%
  \noindent
  \begin{tabularx}{\textwidth}{@{}>{\small\raggedright\arraybackslash}p{2.2cm}@{\hspace{0.4cm}}X@{}}
    #1 & \textbf{#2} \\
    & #3
  \end{tabularx}\vspace{0.3em}
}

% Use blue colored vertical bar via colortbl !-column
\usepackage{colortbl}
\newcolumntype{I}{!{\color{ruleblue}\vrule width 0.5pt}}
\newcommand{\cventryblue}[3]{%
  \noindent
  \begin{tabularx}{\textwidth}{@{}>{\small\raggedright\arraybackslash}p{2.2cm}I@{\hspace{0.4cm}}X@{}}
    #1 & \textbf{#2} \\
    & #3
  \end{tabularx}\vspace{0.3em}
}

% Counters for numbered entries
\newcounter{projectnum}
\newcounter{talknum}
\newcounter{reviewnum}

% Project entry: #1=title, #2=bullet points, #3=label (e.g. Collaborators/Supervisor), #4=names, #5=paper
\newcommand{\projectentry}[5]{%
  \stepcounter{projectnum}
  \vspace{0.5em}\noindent
  {\bfseries \theprojectnum.\ #1}\par\vspace{-0.3em}
  \begin{itemize}[leftmargin=1.2em, itemsep=0pt, parsep=0pt, topsep=0.15em]
    #2
  \end{itemize}\vspace{-0.2em}
  \textit{#3:} #4\par
  \textit{Paper:} #5\par\vspace{0.4em}
}

% Talk entry: #1=conference name, #2=date & location, #3=talk title
\newcommand{\talkentry}[3]{%
  \stepcounter{talknum}
  \vspace{0.4em}\noindent
  {\thetalknum.\ #1}, #2\par\vspace{0.1em}
  \noindent \textit{#3}\par\vspace{0.4em}
}

% Review entry: #1=journal name, #2=paper title
\newcommand{\reviewentry}[2]{%
  \stepcounter{reviewnum}
  \vspace{0.4em}\noindent
  {\thereviewnum.\ \bfseries #1}\par\vspace{0.1em}
  \noindent\textit{#2}\par\vspace{0.4em}
}

% No page numbers
\pagestyle{empty}

% ---------- Document ----------
\begin{document}

% ---------- Title ----------
\begin{center}
  {\large\bfseries Curriculum Vitae}\\[0.3em]
  {\LARGE Abhishek Mishra}
\end{center}

\vspace{0.3em}

% ---------- Contact Info ----------
\noindent
\begin{tabularx}{\textwidth}{@{}X r@{}}
  PhD researcher, & Email: \href{mailto:abhishek.mishra@ulb.be}{abhishek.mishra@ulb.be} \\
  Laboratoire d'Information Quantique, & \href{mailto:abhisheksoloabh@gmail.com}{abhisheksoloabh@gmail.com} \\
  Université Libre de Bruxelles, Belgium. & Website: \href{https://abh1mishra.github.io}{\textbf{https://abh1mishra.github.io}} \\
\end{tabularx}

% ---------- Education ----------
\section*{Education}
\cventryblue{Oct 2022 -- Cont.}{Université Libre de Bruxelles}{%
  PhD in Quantum Information\newline
  Supervisor: Prof.\ Stefano Pironio
}

\cventryblue{Aug 2017 -- Jun 2022}{Indian Institute of Science Education and Research, Kolkata}{%
  BS-MS Dual Degree\,, MS:\,\textnormal{8.9/10 CGPA} \newline
  Thesis title: \textit{QKD security proofs and Device-Independent Certification.}\newline
  Thesis supervisors: Prof.\ Guruprasad Kar \& Dr. Ramij Rahaman, ISI Kolkata.
}

% \cventryblue{2015 -- 2017}{D.A.V Public School Chandrasekharpur, Bhubaneshwar}{%
%   Higher Secondary,\ \ Science(Physics,Chemistry,Maths): 95\%
% }

% ---------- Research Interests ----------
\section*{Research Interests}

\begin{center}
Non-commutative Polynomial Optimization, Quantum Cryptography, Non-commutative Algebra
\end{center}

% ---------- Skills ----------
\section*{Skills}
Julia, NumPy, SciPy, Python, C, Java, Origin, Probability Theory, Linear Algebra, Quantum Information, Convex Optimization, Semi-definite Programming,
Mathematica, App Development, Blockchain.


% ---------- QIP Projects ----------
\section*{Scientific Projects \& Publications}
\setcounter{projectnum}{0}
\projectentry{Partially-commutative Polynomial Optimization}{
  \item Introduced a unified framework interpolating between commutative and non-commutative polynomial optimization, using Gr\"obner basis theory to significantly reduce the complexity of the optimization.
  \item Established foundations of partially-commutative Gr\"obner bases and derived explicit bases for a broad class of problems arising in quantum information.
  \item Developed efficient normal-form algorithm for trace polynomial optimization, superseding prior heuristic approaches.
}{Collaborators}{Dr. Moisés Bermejo Morán, Dr. Stefano Pironio}{To be published soon}

\projectentry{PCPOP : A Julia library for polynomial optimization}{\item Implements normal-form algorithms introduced prior. User-friendly and versatile interface to solve a wide range of polynomial optimization problems with support of sparsity, symmetrization and built-in functions for QKD protocols. \item Efficient and general polynomial library support for Julia. \item Benchmarked against top libraries like OSCAR, QuantumNPA, Ncpol2sdpa etc.\ in terms of efficiency and mathematical rigour.}{Collaborator}{Dr. Moisés Bermejo Morán}{To be published soon,  \href{https://github.com/abh1mishra/PCPOP/tree/main}{Library Link}}

\projectentry{Semi-device Independent QKD}{\item Proposed secure QKD protocols using the restricted-information and restricted-distrust framework in prepare-and-measure scenarios; built numerical framework to demonstrate their robustness. \item Investigating efficient QKD protocols under the basis-independence assumption. }{Collaborators}{Dr. Michele Masini, Dr. Armin Tavakoli, Dr. Stefano Pironio, Dr. Chirag Srivastava}{Two papers To be published soon.}

\projectentry{Transmitter-Device-Independent QKD}{\item Proof-of-principle demonstration of 1Sided Device-Independent QKD over 140km. }{Collaborators}{Dr. Qiang Zeng}{Sent to PRL}

% \projectentry{Master Thesis in Quantum Information (Aug 2021--May 2022)}{
%   \item Understanding the security proof of the BB84 Protocol.
% }{Supervisor}{Prof.\ Guruprasad Kar}{To be published soon}

% \projectentry{Independent Study in Quantum Optics (Aug -- Dec 2021)}{
%   \item Studied Quantum Langevin and Master equations for understanding open quantum systems.
%   \item Worked on Universal Quantum Computation using Continuous Variables.
% }{Supervisor}{Prof.\ Chiranjib Mitra -- IISER Kolkata}{--}

\projectentry{Device-Independent Certification (Part of master thesis)}{
  \item Proved robustness of self-testing of genuine multipartite entanglement.
  \item Built a library for solving NPA hierarchy for any Bell scenario.
}{Collaborator}{Dr.\ Ramij Rahaman}{\href{https://journals.aps.org/pra/abstract/10.1103/PhysRevA.110.L010401}{Self-testing of genuine multipartite entangled states without network assistance}}


% ---------- Contributed Talks ----------
\section*{Contributed Talks \& Posters}
\setcounter{talknum}{0}
\talkentry{828. WE-Heraeus-Seminar on Operator Theory and Polynomial Optimization in Quantum Information Theory}{March 2025, Physikzentrum Bad Honnef}{\textbf{Talk : }Partially-commutative Polynomial Optimization}

\talkentry{Quantum Redemption}{Aug 2024, Mehedeby}{\textbf{Talk : }Partially-commutative Polynomial Optimization}

\talkentry{VERIqTAS workshop}{Mar 2026, Lyon [Accepted]}{\textbf{Talk : }Semi-Device-Independent Quantum Key Distribution Secure under Information and Distrust Assumptions}

\talkentry{POP23 -- Future Trends in Polynomial Optimization}{Nov 2023, LAAS-CNRS, Toulouse}{\textbf{Poster : } \href{https://homepages.laas.fr/henrion/pop23/Bermejo.pdf}{Partially commuting Computations}}

\talkentry{QCRYPT 2024}{Sep 2024, Vigo}{\textbf{Poster : } Partially commuting Computations}
\talkentry{IQOQI Summer School}{Sep 2024, Vienna \& Innsbruck}{\textbf{Poster : } Partially commuting Computations}
\talkentry{YQIS 2025}{Oct 2025, ICFO, Barcelona}{\textbf{Poster : } Partially-commutative Polynomial Optimization}

% % ---------- Contributed Posters ----------
% \section*{Contributed Posters}
% \setcounter{talknum}{0}
% \talkentry{POP23 - Future Trends in Polynomial OPtimization}{Nov 2023, LAAS-CNRS, Toulouse}{\href{https://homepages.laas.fr/henrion/pop23/Bermejo.pdf}{Partially commuting Computations}}

% \talkentry{QCRYPT 2024}{Sep 2024, Vigo}{Partially-commutative Polynomial Optimization}

% \talkentry{YQIS 2025}{Oct 2025, ICFO, Barcelona}{Partially-commutative Polynomial Optimization}
% \talkentry{ICQIF 2022, Feb 2022, ISI Kolkata}{Talk title here}

\section*{Peer-review of Journals}
\setcounter{reviewnum}{0}
Reviewer for international journals, including \textbf{Quantum} and  Elsevier's \textbf{Conmputer Physics Communications}.
% \newpage

% % ---------- Internships ----------
% \section*{Internships}

% \cventryblue{Feb--July 2021}{Internship In Quantum Information}{%
%   Prof.\ Guruprasad Kar -- ISI Kolkata\newline
%   A study project where I learned about mixed-state entanglement and
%   quantum-communication, local distinguishability of pure orthogonal states,
%   quantum random access codes and no masking theorem. It then concluded
%   with a summer school on Quantum Foundations.
% }

% \cventryblue{Summer 2019}{Study of Mesoscopic Systems Using Stochastic Thermodynamics}{%
%   Dr.\ Prabhakar Bhimalapuram -- IIIT Hyderabad\newline
%   Learned about the Langevin equation and Markov chain and how to use them
%   to model stochastic systems. Also learned various statistical tests like the KS
%   test to compare the relevant computed plots to empirical data.
% }

% \cventryblue{Summer 2018}{Study of Nano-scale Properties of Gold colloid and Fractals}{%
%   Prof.\ Anushree Roy -- IIT Kharagpur\newline
%   Learned the usage and basics of Raman Spectroscopy and how to analyze
%   microscope images using ImageJ to determine fractal parameters. Developed a
%   program that fits the image to existing protein growth models and analyzed
%   the gold colloid growth pattern in different media for bio-medical applications.
% }

% % ---------- Relevant Courses ----------
% \section*{Relevant Courses}

% \cventryblue{}{Physics}{%
%   \textnormal{Quantum Information, Quantum Mechanics(I,II,III), Atomic \&\ Optical Physics}
% }
% \vspace{-0.5em}
% \cventryblue{}{Mathematics}{%
%   \textnormal{Linear algebra, Group Theory, Analysis(I,II), Mathematical Physics}
% }
\section*{Research visits}
\cventryblue{Nov-Dec 2025}{Lund university, Sweden}{
  Research visit in the group of Dr. Armin Tavakoli.
}
\cventryblue{Jun 2025}{INRIA, Paris-Saclay}{%
  Research visit in the group of Dr. Marc-oliver Renou.
}
\cventryblue{Aug 2022}{IIT Madras, Chennai}{%
  Project associate job in the Center For Quantum Information, Communication And Computing under Dr. Pradeep Sarvepalli.
}
\cventryblue{Feb-Jul 2021}{ISI Kolkata}{Research visit in the group of Prof. Guruprasad Kar and Dr. Ramij Rahaman.}
\cventryblue{Summer 2018}{IIT Kharagpur}{%
  Research visit in the group of Prof. Anushree Roy.
}
\cventryblue{Summer 2019}{IIIT Hyderabad}{Research visit in the group of Dr. Prabhakar Bhimalapuram.}
% ---------- Teaching Assistant-ship ----------
\section*{Teaching Assistant}

\cventryblue{Spring 2022}{Quantum Information Processing}{%
  \textit{Motivation and introduction, Quantum Algorithms and computation}
}

% ---------- Extra-curricular Activities ----------
\section*{Extra-curricular Activities \& Leadership}

\cventryblue{Web-Development}{Developed Websites or Apps for following:}{%
  \textnormal{Inquivesta 2019, Annual fest of IISER Kolkata\newline
  CR and Admin portal of Inquivesta 2020\newline
  IICM 2021 (Part of it)}
}

\cventryblue{Initiatives}{Started the Coding and Designing Club of IISER Kolkata}{%
  \textnormal{Created the charter of the club. Formed the Council of the Office
  Bearers of the Club. Created the social media platform for the club.}
}

\end{document}
